\thispagestyle{plain}
\normalfont
\enlargethispage{2\baselineskip}

\begin{center}
\vspace*{0.15cm}
{\sffamily\bfseries\zihao{2}霍尔木兹海峡封锁对国际原油价格影响的\\[0.25em]动态建模与情景预测\par}
\vspace{1.1em}
{\sffamily\bfseries\zihao{-4} 摘\quad 要}
\end{center}

\vspace{-0.2em}
\zihao{-4}
霍尔木兹海峡为全球原油运输核心通道，其封锁将通过供应中断、地缘风险、战略储备释放、商业库存缓冲、绕道运输及预期修复共同影响国际油价。本文以布伦特原油期货主力合约价格为数据基础，构建传统供需基准--短期冲击模型--中长期调节模型--敏感性分析的完整模型框架。

本文选取 2017 年 9 月--2026 年 5 月油价数据，以 2026 年 3 月 2 日--5 月 5 日为冲突窗口；该区间内最高收盘价为 114.06 美元/桶，最高盘中价为 119.50 美元/桶。基于线性与常弹性需求的反事实基准测算，理论油价压力区间为 278--337 美元/桶，显著高于实际价格，表明政策缓冲、需求收缩、价格传导与预期修复对油价形成关键约束。

短期冲击模型从物理供需与市场定价双维度建模：纳入供应中断、战略储备（SPR）、商业库存、绕道运输、需求收缩与恐慌衰减；同时引入地缘风险溢价、不确定性、制度风险、缓冲确认与预期修复折价，并规避弹性与需求下降双重计值。模型在 46 个交易日内拟合精度为 RMSE=3.37 美元/桶，MAE=2.74 美元/桶，MAPE=2.73\%；加入 Ridge 收益率辅助层后，RMSE 降至 3.16 美元/桶，方向命中率为 71.1\%。稳健性检验表明模型具备结构稳定性。

中长期调节模型结合 EIA、JODI、OPEC 外生数据与 OVX 隐含波动率，给出 60--180 天乐观、中性、悲观情景。中性情景第 180 天油价约 94.47 美元/桶，悲观情景约 115.90 美元/桶，且存在二次跳涨风险。敏感性分析显示，封锁风险衰减速度、地缘风险权重、制度风险强度为核心影响因子。

引入历史价格状态转移与多约束条件的 Logit--Softmax 模型后，第 180 天油价中位数为 97.54 美元/桶，5\%--95\% 区间为 91.27--112.76 美元/桶，突破 120 美元/桶概率为 76.3\%。检验表明，2026 年冲突窗口具备高风险与高波动特征，外部指标仅用于中长期约束与校验，不参与短期拟合。

\vspace{0.35em}
\noindent{\sffamily\bfseries\zihao{-4} 关键词：} 霍尔木兹海峡；原油价格；短期冲击模型；中长期油价调节模型；战略石油储备；敏感性分析
\normalfont\normalsize
